\documentclass[11pt,a4paper]{article}
\usepackage[margin=1in]{geometry}
\usepackage{booktabs}
\usepackage{amsmath}

\title{Question 1: Word Segmentation and POS Tagging\\
\large English (Brown) vs.\ Spanish (UD Spanish-GSD)}
\date{}

\begin{document}
\maketitle
\vspace{-2.5em}

\section*{Setup}

\begin{center}
\begin{tabular}{lrr}
\toprule
& English (Brown) & Spanish (UD-GSD) \\
\midrule
Train sentences / tokens & 45{,}459 / 804{,}554 & 14{,}186 / 330{,}548 \\
Word types & 41{,}974 & 39{,}748 \\
\textbf{Type/token ratio} & \textbf{0.052} & \textbf{0.120} \\
Mean token length & 4.72 chars & 4.80 chars \\
Dev perplexity (trigram word LM) & 1{,}088.6 & 2{,}233.0 \\
OOV rate (test sample) & 2.44\% & 6.66\% \\
Tagset / morphology-aware tagset & 11 / 11 & 16 / 78 \\
\bottomrule
\end{tabular}
\end{center}

\noindent Spanish reaches almost the same vocabulary from a quarter of the
tokens. Mean token length is identical, so ``Spanish words are longer'' explains
nothing below.

\section*{Q1. Where did the two languages differ most?}

\textbf{In segmentation, and in the unknown-word tail - not in tagging.}

\begin{center}
\begin{tabular}{lrrr}
\toprule
Measure & English & Spanish & Gap \\
\midrule
Segmentation, exact sentences & 66.50\% & 38.75\% & \textbf{+27.75 pp} \\
Segmentation, greedy baseline & 66.71\% & 50.20\% & +16.51 pp \\
Tagging, \emph{unknown} words only & 86.31\% & 76.53\% & \textbf{+9.78 pp} \\
End to end (segment + tag) & 92.56\% & 86.75\% & +5.81 pp \\
Segmentation, trigram DP (token F1) & 95.55\% & 91.59\% & +3.96 pp \\
Tagging, gold segmentation & 96.19\% & 93.75\% & +2.44 pp \\
Tagging, \emph{known} words only & 96.43\% & 94.98\% & +1.45 pp \\
\bottomrule
\end{tabular}
\end{center}

\noindent Once the words are separated correctly the gap nearly closes
(+1.45 pp on known words), so the end-to-end gap is mostly inherited from the
segmenter. The cause is sparse statistics, not grammar: Spanish inflection
splits each lemma across many forms, doubling perplexity and nearly tripling the
OOV rate. Typical failure: \texttt{elcielodespejadoesazul} $\rightarrow$
\texttt{el cielo \textbf{de} \textbf{spejado} es azul}, where a very frequent
preposition is available as the prefix of a rarer content word.

\medskip\noindent\emph{Caveat:} Spanish is scored on 16 tags against English's
11, which slightly flatters English.

\section*{Q2. Did agreement-aware tagging help, or add noise?}

\textbf{It helped in Spanish. In English the question cannot be asked:} Brown
carries no \texttt{FEATS}, so the refined tagset is identical to the plain one
(11 tags either way). That is a fact about the corpus, not about English.

\medskip\noindent Accuracy on the refined 78-label set (91.72\%) is \emph{not}
comparable with the plain tagger's 93.75\% - it is a harder decision. Projected
back to coarse tags, both models sit on the identical 16-way decision:

\begin{center}
\begin{tabular}{lrr}
\toprule
Model & Coarse-tag accuracy & Change \\
\midrule
Plain trigram HMM & 93.75\% & - \\
Morphology-aware HMM, projected & \textbf{94.23\%} & +0.48 pp \\
Plain most-frequent-tag & 88.72\% & - \\
Morphology-aware baseline, projected & 89.74\% & +1.02 pp \\
\bottomrule
\end{tabular}
\end{center}

\noindent So refining the tagset \textbf{added no noise}. But the context-free
baseline gained \emph{more} than the HMM, so this table alone does not prove the
gain comes from agreement modelling. The direct evidence is stronger: in the
context (matching determiner, noun) the transition model prefers an agreeing
adjective over a clashing one by \textbf{12-18$\times$}, and of 2{,}625 adjacent
gold-agreeing pairs it reproduced agreement in \textbf{96.46\%} (DET+ADJ 99.14\%,
ADJ+NOUN 98.72\%, DET+NOUN 97.68\%, NOUN+ADJ 97.15\%; weakest NOUN+NOUN 70.83\%).
End to end this falls to 86.06\%, because 305 pairs contain a word the segmenter
never recovered.

\medskip\noindent Both sample strings are fully correct, including the number
flip propagating across all four agreeing words:

\begin{center}
\texttt{la/DET-Fem-Sg casa/NOUN-Fem-Sg roja/ADJ-Fem-Sg es/AUX-Sg grande/ADJ-Sg}\\
\texttt{las/DET-Fem-Pl casas/NOUN-Fem-Pl rojas/ADJ-Fem-Pl son/AUX-Pl grandes/ADJ-Pl}
\end{center}

\section*{Q3. Segmentation-induced vs.\ genuine tagging errors}

An error is \emph{segmentation-induced} if the gold token's character span was
never recovered (no tag decision was made) and \emph{genuine} if the span was
right and the tag wrong.

\begin{center}
\begin{tabular}{lrrrrr}
\toprule
System & Errors & Seg-induced & Genuine & \% from seg. & Tag acc.\ $|$ span \\
\midrule
\multicolumn{6}{l}{\textbf{English}} \\
Greedy + most-frequent & 2{,}152 & 1{,}861 & 291 & \textbf{86.5\%} & 94.19\% \\
DP segment + most-frequent & 661 & 273 & 388 & 41.3\% & 94.12\% \\
DP segment + HMM tag & 511 & 273 & 238 & \textbf{53.4\%} & 96.39\% \\
\multicolumn{6}{l}{\textbf{Spanish}} \\
Greedy + most-frequent & 4{,}583 & 4{,}162 & 421 & \textbf{90.8\%} & 92.53\% \\
DP segment + most-frequent & 1{,}576 & 831 & 745 & 52.7\% & 91.69\% \\
DP segment + HMM tag & 1{,}298 & 831 & 467 & \textbf{64.0\%} & 94.79\% \\
\bottomrule
\end{tabular}
\end{center}

\noindent With a weak segmenter, segmentation is essentially the \emph{only}
error source (86.5\% / 90.8\%); the baseline tags a correctly segmented token at
94.19\% / 92.53\%, so nearly all its failure is inherited. Improving the tagger
\emph{raises} the segmentation share (41.3\% $\rightarrow$ 53.4\%) because it
removes genuine errors and cannot touch inherited ones - that number going up
is the tagger working. Spanish stays segmentation-dominated (64.0\% vs.\ 53.4\%).

\medskip\noindent The two classes look nothing alike:
segmentation-induced - \texttt{unfailing} $\rightarrow$ \texttt{un + failing},
\texttt{dallasbased} $\rightarrow$ \texttt{dallas + based};
genuine - \texttt{eran} VERB$\rightarrow$AUX, \texttt{que}
PRON$\rightarrow$SCONJ. Top confusions: English VERB/NOUN (35+24) and PRT/ADP
(17+16); Spanish PROPN/NOUN (75+50, unsurprising once case is normalised away)
and PRON/DET (42).

\section*{Q4. How much better than the simple baselines?}

\begin{center}
\begin{tabular}{lrrrr}
\toprule
Task & Baseline & Model & Gain & Error reduction \\
\midrule
\multicolumn{5}{l}{\textbf{English}} \\
Segmentation (token F1) & 66.71\% & 95.55\% & +28.85 pp & \textbf{86.6\%} \\
Segmentation (exact sentences) & 16.25\% & 66.50\% & +50.25 pp & 60.0\% \\
Tagging (gold segmentation) & 93.64\% & 96.19\% & +2.55 pp & 40.0\% \\
End to end & 68.68\% & 92.56\% & +23.88 pp & \textbf{76.3\%} \\
\multicolumn{5}{l}{\textbf{Spanish}} \\
Segmentation (token F1) & 50.20\% & 91.59\% & +41.39 pp & \textbf{83.1\%} \\
Segmentation (exact sentences) & 5.75\% & 38.75\% & +33.00 pp & 35.0\% \\
Tagging (gold segmentation) & 88.72\% & 93.75\% & +5.03 pp & 44.6\% \\
End to end & 53.21\% & 86.75\% & +33.54 pp & \textbf{71.7\%} \\
\bottomrule
\end{tabular}
\end{center}

\noindent \textbf{Segmentation is where the model pays for itself}, removing
83-87\% of the baseline's errors: greedy longest-match cannot reconsider a bite,
so one wrong match corrupts every boundary after it (only 16.25\% / 5.75\% of
sentences perfect, vs.\ 66.50\% / 38.75\%). \textbf{Tagging is where the baseline
is already strong} - most word types are unambiguous - so +2.55 / +5.03 pp
looks small but removes 40-45\% of the remaining errors. Both baselines were
given every advantage: greedy gets the \emph{full} training vocabulary including
hapax words the LM itself discards.

\medskip\noindent Cost: greedy runs in 0.10\,ms/sentence against 92\,ms for the
DP decoder (English), i.e.\ three orders of magnitude for 29-41 points of F1.
Beam search makes that affordable - at width 8 it matches exact DP to four
significant figures (96.53\% vs.\ 96.53\%; 92.91\% vs.\ 92.91\%) while running
10-20$\times$ faster.

\section*{Known failures}

\begin{itemize}\itemsep2pt
\item \texttt{brown} $\rightarrow$ NOUN, not ADJ: in the Brown corpus the form is
predominantly a proper noun.
\item \texttt{despejado} $\rightarrow$ \texttt{de + spejado} (Q1).
\item \texttt{pueden} $\rightarrow$ AUX where the brief writes VERB: a UD tagset
convention, not an error.
\item English morphology was never tested - Brown has no \texttt{FEATS}.
\item All test figures are 400 sentences per language; differences below roughly
half a point should not be read as real.
\end{itemize}

\end{document}
